\documentclass{article}
\usepackage{amsthm}
\usepackage{mathpartir}
\usepackage{url}
\usepackage{imprimatur}
\newtheorem{theorem}{Theorem}
% The paper's environments, as imprimatur.toml's [aliases] declares them.
\leanblockalias{leangrammar}{}
\leanblockalias{leanrules}{}
\leanblockalias{leantheorem}{theorem}
\leanblockalias{leanconventions}{}
\leanblockalias{leannotation}{}
% Conventions of the generated rules: \val is a value, and \rulelabel{Key}{label}
% makes a rule a target for \ref.
\newcommand\val{v}
\makeatletter
\newcommand\rulelabel[2]{\def\@currentlabel{#2}\label{rule:#1}}
\makeatother
\begin{document}

\section{A small calculus, generated from Lean}

This paper presents a simply-typed $\lambda$-calculus whose definitions and
results are in Lean. Everything in a \texttt{lean...} environment below is
generated from the Lean development by imprimatur; the rest is written by
hand, and the hand-written statements are checked against Lean.

\subsection{Syntax}

Terms are name-carrying: a variable is a name, and an abstraction binds a
name. The grammar is generated from the inductive types. The descriptions on
the right are the authors': imprimatur keeps them when it regenerates the
grammar, and uses the constructors' docstrings only for new ones.

\begin{leangrammar}{Stlc.Ty}{Types}
\end{leangrammar}

\begin{leangrammar}{Stlc.Tm}{Terms}
\end{leangrammar}

\subsection{Names, $\alpha$-equivalence and substitution}

We handle names as Pitts does~\cite{Pitts06}. The basic operation is
\emph{swapping}: $(a\ b) \cdot e$ exchanges the names $a$ and $b$
everywhere in $e$, bound, binding and free occurrences alike. Unlike
renaming, swapping needs no side condition at a binder, and it is its own
inverse; its definition is structural:

\begin{center}
\begin{leanblock}[view=math, def]{Stlc.Tm.swap}{Swapping}
\end{leanblock}
\end{center}

Two abstractions are $\alpha$-equivalent when swapping each binder's name
with a name $c$ fresh for both, written $c \mathrel{\#} \ldots$, makes their
bodies $\alpha$-equivalent:

\begin{leanrules}{Stlc.Alpha}{Alpha-equivalence}
\end{leanrules}

Substitution $e[u/x]$ is then defined by recursion on $e$. At an abstraction
$\lambda a{:}\tau.\, e$ it first swaps $a$ with a name $c$ fresh for $x$,
$a$, $e$ and $u$, so that no free name of $u$ is captured, and continues in
$(a\ c) \cdot e$\leancheck{Stlc.subst}. Swapping preserves the size of a
term\leancheck{Stlc.Tm.size_swap}, so the recursion terminates, although
$(a\ c) \cdot e$ is not a subterm of the abstraction.

\subsection{Typing}

The typing rules are generated from the constructors of \texttt{HasType}.
Their labels come from the constructors' docstrings. The well-formedness side
conditions of the Lean rules are left implicit, by a convention declared in
the notation module.

\begin{leanrules}{Stlc.HasType}{Typing}
\end{leanrules}

\subsection{Reduction}

Reduction is call by value. A premise that a term is a value is left
implicit: the value is written $\val$ instead. Each rule is labelled from its
constructor's name, and can be referred to: rule~\ref{rule:AppLeft}
reduces the function of an application.

\begin{leanrules}{Stlc.Step}{Reduction}
\end{leanrules}

\subsection{Results}

A result can be generated, as a sentence; its key can be a title, declared
in the notation module:

\begin{leantheorem}{Swapping is involutive}{Swapping is involutive}
\end{leantheorem}

\begin{leantheorem}{Values are final}{Values are final}
\end{leantheorem}

Or as a rule, from a theorem whose hypotheses become its premises:
\begin{center}
\begin{leanblock}[view=math, rule]{Stlc.step_app}{Congruence}
\end{leanblock}
\end{center}

Or the prose can be written by hand and checked: the unit value has only
the base type, that is, if $\Gamma \vdash () : \tau$ then
$\tau = \mathsf{b}$\leancheck{Stlc.unit_type}. Should the Lean statement of
\verb|unit_type| change, \texttt{imprimatur check} reports this sentence.

\subsection{Type safety}

The substitution lemma is proved as Pitts's paper suggests. Typing is
equivariant: swapping two names throughout a term and its context preserves
its type.

\begin{leantheorem}{Equivariance}{Equivariance}
\end{leantheorem}

Typing also depends on the context only at the term's free names, which
gives weakening and shadowing at once\leancheck{Stlc.HasType.agree}. So
substitution preserves types: if $\Gamma, x{:}\tau_1 \vdash e : \tau$ and
$\Gamma \vdash u : \tau_1$, then $\Gamma \vdash e[u/x] : \tau$. At an
abstraction, substitution renames the binder to a fresh name by swapping, and
equivariance carries the typing of its body across the renaming; the
induction is on the size of $e$, which swapping
preserves\leancheck{Stlc.subst_typing}. Type safety follows.

\begin{leantheorem}{Preservation}{Preservation}
\end{leantheorem}

\begin{leantheorem}{Progress}{Progress}
\end{leantheorem}

\begin{thebibliography}{1}
\bibitem[Pitts(2006)]{Pitts06}
Andrew~M. Pitts.
\newblock Alpha-structural recursion and induction.
\newblock \emph{Journal of the ACM}, 53(3):459--506, 2006.
\newblock \url{https://doi.org/10.1145/1147954.1147961}.
\end{thebibliography}

\appendix
\section{Index of notation}\label{app:notation}

Metavariables range over the types below; a variable with a subscript or a
prime ranges over the same type. imprimatur generates the table from the
categories the notation declares, and checks that every variable printed with
a metavariable has its type.

\begin{leanconventions}{Conventions}{Metavariables and the types they range over}
\end{leanconventions}

The notation, grouped by where it is defined: the paper's own (\textit{Stlc})
and Lean's. Each is written as the paper writes it, with the Lean constant it
stands for.

\begin{leannotation}{Notation}{The notation of the paper}
\end{leannotation}

\end{document}
